\documentclass{beamer}
\usetheme{Madrid}
\usecolortheme{default}
\usepackage{amsmath}
\usepackage{booktabs}

% Add this line to fix the slide counter
\usepackage{appendixnumberbeamer}

\setbeamertemplate{navigation symbols}{}

\title[Electoral Accountability]{Electoral Accountability and Economic Policy}
\subtitle{Based on Besley \& Case (1995)}
\author{Lucas Sernik}
\date{\today}

\begin{document}

\begin{frame}
\titlepage
\end{frame}

%--- MOTIVATION ---%
\begin{frame}{Motivation}
\begin{itemize}
    \item A central question in political economy: \textbf{do electoral incentives discipline politicians?}
    \vspace{0.3em}
    \item Do politicians act in the public interest, or do they use political power for personal gain?
    \vspace{0.3em}
    \item The desire to maintain a reputation is often thought to keep politicians in check
    \vspace{0.3em}
    \item Officials who care to run again must act in voters' interest to merit reelection
    \vspace{0.3em}
    \item But what happens when reelection is no longer possible?
    \vspace{0.3em}
    \item \textbf{Besley \& Case (1995)} exploit gubernatorial term limits in U.S. states to test this
\end{itemize}
\end{frame}

%--- RESEARCH QUESTION ---%
\begin{frame}{Research Question}
\begin{block}{Central Question}
Does the removal of reelection incentives --- through a binding term limit --- alter economic policy choices?
\end{block}
\vspace{1em}
\begin{itemize}
    \item If incumbents care about reputation, they should behave differently when they cannot run again
    \vspace{0.5em}
    \item A binding term limit removes the future electoral reward that keeps politicians in check, allowing us to observe what happens when accountability disappears.
    \vspace{0.5em}
    \item This allows a clean test of the \textbf{political agency model}
\end{itemize}
\end{frame}

%--- THEORETICAL FRAMEWORK ---%
\begin{frame}{Theoretical Framework: Reputation-Building Model}
\begin{itemize}
    \item Based on Banks \& Sundaram (1993)
    \vspace{0.5em}
    \item Each governor has an unobservable \textbf{type} $\omega \in \{\omega_1, ..., \omega_N\}$
    \vspace{0.5em}
    \item Takes unobservable \textbf{effort} $\alpha$, which affects policy outcome $r$
    \vspace{0.5em}
    \item Voters observe $r$ and use it to infer the governor's type via Bayes' rule
\end{itemize}
\vspace{0.5em}
\begin{block}{Two Regimes}
\begin{enumerate}
    \item \textbf{Term-limited (lame duck)}: maximizes immediate payoffs --- puts in less effort
    \item \textbf{Reelection-eligible}: puts in extra effort to build reputation and secure reelection
\end{enumerate}
\end{block}
\end{frame}

%--- MODEL ---%
\begin{frame}{The Model}
\textbf{Short-run (term-limited) decision:}
\begin{equation*}
    \alpha_s(\omega) = \arg\max_{\alpha} \left\{ u(\alpha, \omega) \; | \; \alpha \in [\underline{\alpha}, \bar{\alpha}] \right\}
\end{equation*}
\vspace{0.5em}
\textbf{Long-run (reelection-eligible) decision:}
\begin{equation*}
    \alpha_l(\omega) = \arg\max_{\alpha} \left\{ u(\alpha, \omega) + \delta \Pr\left[r \in R(\sigma) \; | \; \alpha\right] u(\alpha_s(\omega), \omega) \right\}
\end{equation*}
\vspace{0.5em}
\begin{block}{Key Proposition}
Incumbents eligible for reelection put in \textbf{more effort} than term-limited ones. Those in their last term give lower payoffs to voters, on average.
\end{block}
\end{frame}

%--- EMPIRICAL STRATEGY ---%
\begin{frame}{Empirical Strategy}
\textbf{Estimated equation:}
\begin{equation*}
    P_{st} = \zeta_s + \psi_t + \gamma T_{st} + \alpha Z_{st} + \varepsilon_{st}
\end{equation*}
where:
\begin{itemize}
    \item $P_{st}$: policy outcome (taxes, spending, minimum wage) in state $s$, year $t$
    \item $\zeta_s$: state fixed effects
    \item $\psi_t$: year fixed effects
    \item $T_{st} = 1$ if governor cannot run again (binding term limit)
    \item $Z_{st}$: controls (income per capita, demographics, population)
    \item $\gamma$: \textbf{coefficient of interest}
\end{itemize}
\vspace{0.3em}
\begin{alertblock}{}
If $\gamma = 0$, the reputation-building model is not supported.
\end{alertblock}
\end{frame}

%--- RESULTS ---%
\begin{frame}{Main Results}
\begin{table}
\centering
\small
\begin{tabular}{lcc}
\toprule
\textbf{Outcome} & \textbf{Coefficient on $T_{st}$} & \textbf{Significance} \\
\midrule
Sales taxes per capita      & $+7.86$  & ** \\
Income taxes per capita     & $+8.74$  & ** \\
Corporate taxes per capita  & $+0.57$  & --- \\
Total taxes per capita      & $+6.71$  & --- \\
State expenditure per capita & $+14.38$ & ** \\
State minimum wage          & $-0.14$  & ** \\
\bottomrule
\multicolumn{3}{l}{\footnotesize ** $p < 0.05$. All variables per capita in 1982 dollars.}
\end{tabular}
\end{table}
\vspace{0.5em}
\begin{itemize}
    \item Term-limited governors raise taxes and spending
    \item They reduce minimum wages --- consistent with less resistance to special interests
\end{itemize}
\end{frame}

%--- PARTY HETEROGENEITY ---%
\begin{frame}{Party Heterogeneity}
\begin{block}{Key Finding}
The term-limit effect is \textbf{driven entirely by Democratic governors}.
\end{block}
\vspace{0.5em}
\begin{itemize}
    \item Democratic lame ducks: raise sales, income, and corporate taxes significantly
    \item Republican lame ducks: no significant effect on taxes or spending
    \item Suggests \textbf{incomplete party discipline} --- individual incentives matter
    \vspace{0.5em}
    \item Fiscal cycle interpretation: Democrats hold spending down in first term to get reelected, then increase it in the lame-duck term
\end{itemize}
\end{frame}

%--- FISCAL CYCLE ---%
\begin{frame}{The Fiscal Cycle}
\begin{itemize}
    \item The main effect of term limits is not a permanent shift but a \textbf{fiscal cycle}
    \vspace{0.5em}
    \item Governors hold spending \textbf{below} the state mean in their first term
    \item Governors hold spending \textbf{above} the state mean in their lame-duck term
    \vspace{0.5em}
    \item No significant within-term variation in behavior --- the shift is sharp at the term boundary
    \vspace{0.5em}
    \item Effect on \textbf{state income per capita}: lame-duck Democratic governors are associated with a 1--2\% reduction in state income --- the efficiency cost of the fiscal cycle
\end{itemize}
\end{frame}

%--- RELEVANCE FOR OUR PAPER ---%
\begin{frame}{Relevance for Our Paper}
\begin{block}{Besley \& Case (1995) establish:}
\begin{enumerate}
    \item Reelection incentives discipline politicians' fiscal behavior
    \item Term limits provide exogenous variation in electoral incentives
    \item The effect operates through reputation-building, not party discipline
\end{enumerate}
\end{block}
\vspace{0.5em}
\begin{alertblock}{Our contribution}
We apply this framework to \textbf{Brazilian municipalities} in São Paulo, asking whether reelection eligibility affects \textbf{spending composition} --- not just the level of taxes and spending --- and whether this effect is moderated by \textbf{fiscal discretion}.
\end{alertblock}
\end{frame}

%--- CONTEXT: BRAZILIAN MUNICIPAL ELECTIONS ---%
\begin{frame}{Brazilian Municipal Elections}
\begin{itemize}
    \item Brazilian municipalities hold \textbf{direct elections} for mayor every four years
    \vspace{0.5em}
    \item Elections are held simultaneously across all municipalities --- in our sample, the relevant elections occurred in \textbf{2016} and \textbf{2020}
    \vspace{0.5em}
    \item Voter turnout is high by international standards, partly due to \textbf{compulsory voting} for citizens aged 18--70
    \vspace{0.5em}
    \item São Paulo state alone has \textbf{645 municipalities}, ranging from the capital with over 12 million inhabitants to small rural towns with fewer than 1,000 residents
\end{itemize}
\end{frame}

%--- TERM LIMITS ---%
\begin{frame}{Term Limits in Brazilian Municipalities}
\begin{itemize}
    \item Brazilian law limits mayors to \textbf{two consecutive terms} in office
    \vspace{0.5em}
    \item A mayor elected in 2016 for the \textbf{first time} is eligible to run again in 2020
    \vspace{0.5em}
    \item A mayor elected in 2016 who was \textbf{already serving their first term} since 2012 faces a binding term limit and cannot run in 2020
    \vspace{0.5em}
    \item This generates clean variation in reelection incentives \textbf{within the same election cycle} across municipalities
\end{itemize}
\begin{block}{Identification}
    Comparing mayors who can run again in 2020 against those who cannot, conditional on controls, isolates the causal effect of reelection incentives on spending behavior.
\end{block}
\end{frame}

%--- INSTITUTIONAL SETTING: WHO CONTROLS SPENDING? ---%
\begin{frame}{Who Controls Municipal Spending?}
\begin{itemize}
    \item Brazilian municipalities operate under a \textbf{separation of powers}: the \textit{Câmara dos Vereadores} (city council) approves laws and the annual budget, while the \textbf{mayor executes it}
    \vspace{0.5em}
    \item In practice, the mayor has \textbf{substantial discretion} over spending composition through several channels:
    \vspace{0.5em}
    \begin{enumerate}
        \item \textbf{Budget proposal}: the mayor drafts the annual budget (\textit{Lei Orçamentária Anual}), which the council typically approves with minor modifications
        \vspace{0.3em}
        \item \textbf{Execution discretion}: the mayor controls the \textit{timing} and \textit{allocation} of expenditures within approved categories
        \vspace{0.3em}
        \item \textbf{Supplementary credits}: the mayor can open \textit{créditos suplementares} to reallocate spending across categories without full council approval
        \vspace{0.3em}
        \item \textbf{Administrative decisions}: hiring, procurement, and contracting decisions --- which shape spending composition in practice --- are executive prerogatives
    \end{enumerate}
\end{itemize}
\end{frame}

%--- FISCAL DISCRETION ---%
\begin{frame}{The Limits of Mayoral Discretion}
\begin{itemize}
    \item Not all spending is equally subject to mayoral discretion
    \vspace{0.5em}
    \item A large share of municipal revenue is \textbf{constitutionally earmarked}:
    \vspace{0.5em}
    \begin{itemize}
        \item \textbf{FUNDEB}: at least 25\% of tax revenue must be spent on education
        \item \textbf{SUS transfers}: federal health transfers are earmarked exclusively to health spending
        \item \textbf{Convênios}: federal and state grants are tied to specific programs
    \end{itemize}
    \vspace{0.5em}
    \item What remains after earmarked transfers is \textbf{treasury revenue} (\textit{tesouro}) --- the discretionary budget the mayor truly controls
    \vspace{0.5em}
    \item This motivates our key moderating variable: the \textbf{share of total spending funded by discretionary treasury resources}, $D_{it}$, which measures how much fiscal power the mayor actually has
\end{itemize}
\end{frame}

%--- RESEARCH QUESTION ---%
\begin{frame}{Our Research Question}
\begin{block}{Central Question}
Do reelection incentives affect how Brazilian mayors allocate their \textbf{discretionary} budget across spending categories?
\end{block}
\vspace{1em}
\begin{itemize}
    \item We follow Besley \& Case (1995) in using \textbf{term limits} as a source of exogenous variation in reelection incentives
    \vspace{0.5em}
    \item Our contribution is to focus on \textbf{spending composition} rather than aggregate levels, and to explicitly account for \textbf{fiscal discretion} as a moderating force
    \vspace{0.5em}
    \item The Brazilian context is particularly informative because \textbf{heavy earmarking} limits mayoral discretion --- if political budget cycles exist, they should be concentrated in the discretionary budget
\end{itemize}
\end{frame}

%--- DATA SOURCES ---%
\begin{frame}{Data Sources}
\begin{itemize}
    \item \textbf{Municipal expenditure data}: Tribunal de Contas do Estado de São Paulo (TCESP)
    \begin{itemize}
        \item Annual expenditure records by spending category and funding source
        \item 643 municipalities, 2017--2024
    \end{itemize}
    \vspace{0.5em}
    \item \textbf{Electoral data}: Tribunal Superior Eleitoral (TSE)
    \begin{itemize}
        \item Candidate-level data for 2016 and 2020 municipal elections
        \item Reelection eligibility, vote shares, number of candidates
    \end{itemize}
    \vspace{0.5em}
    \item \textbf{Socioeconomic controls}: IBGE Census 2010 and Atlas Brasil
    \begin{itemize}
        \item Population, PIB per capita, age distribution
        \item Municipal HDI, sanitation access, infant mortality
    \end{itemize}
    \item All socioeconomic controls are measured in 2010, eliminating reverse causality concerns.
\end{itemize}
\end{frame}

%--- DEPENDENT VARIABLE ---%
\begin{frame}{Dependent Variable: Discretionary Spending Share}
\begin{itemize}
    \item Standard measures of municipal spending conflate \textbf{mayoral choices} with \textbf{mechanically earmarked transfers}
    \vspace{0.5em}
    \item We focus instead on the \textbf{discretionary spending share}:
\end{itemize}
\vspace{0.5em}
\begin{block}{Dependent Variable}
\begin{equation*}
    s^{d}_{ikt} = \frac{\text{Discretionary spending in category } k_{it}}{\text{Total discretionary spending}_{it}} \times 100
\end{equation*}
\end{block}
\vspace{0.5em}
\begin{itemize}
    \item \textbf{Discretionary spending} = expenditure funded by \textit{tesouro} or \textit{recursos próprios da administração indireta}
    \item \textbf{Earmarked spending} = federal/state transfers and special funds (\textit{FUNDEB}, \textit{SUS}, \textit{convênios})
    \vspace{0.5em}
    \item $s^{d}_{ikt}$ measures \textbf{how much of total discretionary money did the mayor choose to put in a given category} --- a genuine revealed preference for spending priority.
\end{itemize}
\end{frame}

%--- KEY VARIABLES ---%
\begin{frame}{Key Variables}
\begin{itemize}
    \item \textbf{Treatment variable} $R_{i_t} \in \{0,1\}$:
    \item $R_{i_t} \in \{0,1\}$ varies at the \textbf{municipality-term} level:
\begin{itemize}
    \item For 2017--2020: $R_{i_1} = 1$ if the 2016 winner was in their \textbf{first term}, eligible to run in 2020; $R_{i_1} = 0$ if they had already served since 2012
    \item For 2021--2024: $R_{i_2} = 1$ if the 2020 winner was in their \textbf{first term}, eligible to run in 2024; $R_{i_2} = 0$ if they had already served since 2016
\end{itemize}
\item Predetermined relative to spending: determined by prior election outcomes before each term begins
    \vspace{0.5em}
    \item \textbf{Moderating variable} $D_{it}$:
    \begin{itemize}
        \item Discretionary share of total expenditure:
        \begin{equation*}
            D_{it} = \frac{\text{Total discretionary spending}_{it}}{\text{Total spending}_{it}} \times 100
        \end{equation*}
        \item Measures \textbf{mayoral fiscal power} --- how much of the total budget the mayor truly controls
        \item Varies at the municipality-year level as federal transfer rules and revenues change
    \end{itemize}
\end{itemize}
\end{frame}

%--- REGRESSION EQUATION ---%
\begin{frame}{Regression Equation}
\begin{block}{Estimated Model}
\begin{equation*}
    s^{d}_{ikt} = \beta_0
                + \beta_1 R_{it}
                + \beta_2 \tilde{D}_{it}
                + \beta_3 (R_{it} \times \tilde{D}_{it})
                + \beta_4 E_{it}
                + \beta_5 N_{it}
                + \mathbf{X}_i'\boldsymbol{\gamma}
                + \mu_t
                + \varepsilon_{it}
\end{equation*}
\end{block}
\vspace{0.3em}
\begin{itemize}
    \item $s^{d}_{ikt}$: discretionary spending share in category $k$, municipality $i$, year $t$
    \item $R_{it}$: reelection eligibility --- varies at municipality-term level
    \item $\tilde{D}_{it}$: demeaned discretionary share of total expenditure
    \item $R_{it} \times \tilde{D}_{it}$: interaction --- does greater fiscal discretion amplify the effect of reelection incentives on spending composition?
    \item $E_{it}$: total expenditure per capita
    \item $N_{it}$: number of candidates in the relevant election --- varies at municipality-term level
    \item $\mathbf{X}_i$: 2010 Census controls (PIB per capita, IDHM, share of children)
    \item $\mu_t$: year fixed effects
\end{itemize}
\end{frame}

%--- COEFFICIENTS OF INTEREST ---%
\begin{frame}{Coefficients of Interest}
\begin{itemize}
    \item $\beta_1$ --- \textbf{main effect of reelection eligibility}:
    \begin{itemize}
        \item Effect of reelection eligibility on discretionary spending composition \textbf{at the sample mean of fiscal discretion} ($\bar{D} = 71.4\%$)
        \item If $\beta_1 \neq 0$: reelection incentives distort discretionary spending composition
        \item If $\beta_1 = 0$: no political budget cycle in discretionary spending
    \end{itemize}
    \vspace{0.5em}
    \item $\beta_2$ --- \textbf{effect of fiscal discretion on spending composition}:
    \begin{itemize}
        \item How spending shares change as fiscal discretion moves \textbf{above or below the sample mean} of 71.4\%
        \item Positive $\beta_2$: above average fiscal discretion $\rightarrow$ more spending in that category
        \item Negative $\beta_2$: above average fiscal discretion $\rightarrow$ less spending in that category
    \end{itemize}
\end{itemize}
\end{frame}

\begin{frame}{Coefficients of Interest}
\begin{itemize}
    \item $\beta_3$ --- \textbf{moderation effect of fiscal discretion on reelection incentives}:
    \begin{itemize}
        \item How the spending gap between reelection-eligible and term-limited mayors changes as fiscal discretion moves above or below its sample mean
        \item Negative $\beta_3$: gap \textbf{narrows} for above-average discretion municipalities
        \item Positive $\beta_3$: gap \textbf{widens} for above-average discretion municipalities
    \end{itemize}
\end{itemize}
\end{frame}

%--- SPENDING CATEGORIES ---%
\begin{frame}{Spending Categories}
\begin{itemize}
    \item We estimate the model separately for \textbf{six functional spending categories} from TCESP data
\end{itemize}
\vspace{0.5em}
\begin{table}
\centering
\small
\begin{tabular}{lp{7cm}}
\toprule
\textbf{Category} & \textbf{Description} \\
\midrule
Health (\textit{Saúde}) & Primary healthcare, hospitals, and public health programs \\
Education (\textit{Educação}) & Schools, teachers, and educational infrastructure \\
Urbanismo & Urban infrastructure, roads, and public works \\
Social Assistance & Welfare programs and social protection \\
Administration & General government overhead and bureaucracy \\
Transport & Public transportation and road maintenance \\
\bottomrule
\end{tabular}
\end{table}
\vspace{0.5em}
\begin{itemize}
    \item Health and education are heavily constrained by \textbf{constitutional earmarking} (SUS, FUNDEB) --- the reelection effect and its moderation by fiscal discretion are expected to be concentrated in the remaining categories where mayors have genuine room to maneuver
\end{itemize}
\end{frame}

%--- RESULTS OVERVIEW ---%
\begin{frame}{Results Overview}
\begin{itemize}
    \item We estimate the model separately for six spending categories
    \vspace{0.5em}
    \item Results are organized around three findings:
    \vspace{0.5em}
    \begin{enumerate}
        \item \textbf{Reelection eligibility} ($\beta_1$) significantly affects discretionary spending in \textbf{urbanismo} and \textbf{administration}
        \vspace{0.3em}
        \item \textbf{Fiscal discretion} ($\beta_2$) shapes spending composition independently of electoral incentives in \textbf{social assistance}, \textbf{transport}, and \textbf{urbanismo}
        \vspace{0.3em}
        \item The \textbf{moderation effect} ($\beta_3$) is statistically significant in \textbf{social assistance} and \textbf{administration}, but economically negligible
    \end{enumerate}
    \vspace{0.5em}
    \item Health, education, and transport show \textbf{no significant reelection effect} --- consistent with constitutional earmarking limiting mayoral discretion in these categories
\end{itemize}
\end{frame}

%--- FINDING 1: SUMMARY TABLE ---%
\begin{frame}{Finding 1: Reelection Eligibility and Spending Composition}
\begin{table}
\centering
\resizebox{\textwidth}{!}{%
\begin{tabular}{lcccc}
\hline\hline
 & \multicolumn{4}{c}{Coefficient on $R_{it}$ ($\beta_1$)} \\
\cmidrule(lr){2-5}
Category & Estimate & Std. Error & T-stat & P-value \\
\hline
Health          & $-$0.093 & 0.279 & $-$0.34 & 0.738 \\
Education       & 0.065    & 0.199 & 0.32    & 0.746 \\
Urbanismo       & $-$0.963$^{***}$ & 0.299 & $-$3.22 & 0.001 \\
Social Assist.  & 0.080    & 0.100 & 0.80    & 0.425 \\
Administration  & 1.027$^{***}$    & 0.373 & 2.75    & 0.006 \\
Transport       & 0.023    & 0.201 & 0.11    & 0.910 \\
\hline\hline
\multicolumn{5}{l}{\footnotesize $^{***}p<0.01$, $^{**}p<0.05$, $^{*}p<0.10$. Clustered SE at municipality level.}
\end{tabular}%
}
\end{table}
\vspace{0.5em}
\begin{itemize}
    \item Reelection eligibility affects spending composition \textbf{only through urbanismo and administration}
\end{itemize}
\end{frame}

%--- FINDING 1A: URBANISMO SLIDE 1 ---%
\begin{frame}{Finding 1: Reelection Eligibility and Urbanismo}
\begin{block}{Result}
Reelection-eligible mayors allocate \textbf{0.96 percentage points less} of their discretionary budget to urbanismo relative to term-limited mayors ($\beta_1 = -0.963$, $p = 0.001$)
\end{block}
\vspace{1em}
\begin{itemize}
    \item Urbanismo covers \textbf{urban infrastructure, roads, and public works} --- a category historically associated with corruption in Brazilian municipalities
    \vspace{1em}
    \item We interpret the \textbf{negative} $\beta_1$ as evidence of \textbf{electoral accountability constraining opportunistic spending}: reelection-eligible mayors restrain infrastructure spending because corruption in public works carries electoral risk before the next election
\end{itemize}
\end{frame}

%--- FINDING 1A: URBANISMO SLIDE 2 ---%
\begin{frame}{Finding 1: Reelection Eligibility and Urbanismo}
\begin{block}{Result}
Reelection-eligible mayors allocate \textbf{0.96 percentage points less} of their discretionary budget to urbanismo relative to term-limited mayors ($\beta_1 = -0.963$, $p = 0.001$)
\end{block}
\vspace{1em}
\begin{itemize}
    \item Term-limited mayors, facing \textbf{no future electoral punishment}, increase discretionary urbanismo spending --- consistent with rent extraction through inflated construction contracts, a well-documented phenomenon in Brazilian public administration
    \vspace{1em}
    \item This is consistent with the \textbf{reputation-building model}: accountability disciplines mayors when reelection is at stake, and opportunistic behavior emerges when it is not
\end{itemize}
\end{frame}

%--- FINDING 1B: ADMINISTRATION SLIDE 1 ---%
\begin{frame}{Finding 1: Reelection Eligibility and Administration}
\begin{block}{Result}
Reelection-eligible mayors allocate \textbf{1.03 percentage points more} of their discretionary budget to administration relative to term-limited mayors ($\beta_1 = 1.027$, $p = 0.006$)
\end{block}
\vspace{1em}
\begin{itemize}
    \item Administration covers \textbf{general government overhead}, including public hiring and bureaucratic expenditure
    \vspace{1em}
    \item We interpret this as evidence of \textbf{political patronage}: reelection-seeking mayors strategically expand the municipal payroll to generate electoral loyalty
\end{itemize}
\end{frame}

%--- FINDING 1B: ADMINISTRATION SLIDE 2 ---%
\begin{frame}{Finding 1: Reelection Eligibility and Administration}
\begin{block}{Result}
Reelection-eligible mayors allocate \textbf{1.03 percentage points more} of their discretionary budget to administration relative to term-limited mayors ($\beta_1 = 1.027$, $p = 0.006$)
\end{block}
\vspace{1em}
\begin{itemize}
    \item This is consistent with the broader literature on \textbf{political clientelism} in Brazilian municipalities, where public employment is a well-documented tool of electoral mobilization
    \vspace{1em}
    \item Importantly, this effect is identified \textbf{within the discretionary budget} --- it reflects a genuine mayoral choice to prioritize administrative hiring over service delivery when reelection is at stake
\end{itemize}
\end{frame}

%--- FINDING 2A: FISCAL DISCRETION SLIDE 1 ---%
\begin{frame}{Finding 2: Fiscal Discretion and Spending Composition}
\begin{block}{Result}
The discretionary share of total expenditure ($\tilde{D}_{it}$) significantly shapes spending composition in three categories, independently of electoral incentives
\end{block}
\vspace{0.5em}
\begin{table}
\centering
\resizebox{\textwidth}{!}{%
\begin{tabular}{lcccc}
\hline\hline
Category & Estimate & Std. Error & T-stat & P-value \\
\hline
Social Assist.  & $+$0.054$^{***}$ & 0.015 & 3.63 & 0.000 \\
Transport       & $+$0.069$^{***}$ & 0.025 & 2.74 & 0.006 \\
Urbanismo       & $-$0.095$^{**}$  & 0.041 & $-$2.33 & 0.020 \\
\hline\hline
\multicolumn{5}{l}{\footnotesize $^{***}p<0.01$, $^{**}p<0.05$, $^{*}p<0.10$.}
\end{tabular}%
}
\end{table}
\vspace{0.5em}
\end{frame}

\begin{frame}{Finding 2: Fiscal Discretion and Spending Composition}
\begin{block}{Result}
The discretionary share of total expenditure ($\tilde{D}_{it}$) significantly shapes spending composition in three categories, independently of electoral incentives
\end{block}
\vspace{1em}
\begin{itemize}
    \item Municipalities with \textbf{above average fiscal discretion} allocate proportionally more to social assistance and transport --- when mayors have more budget freedom than the typical municipality, these categories reflect \textbf{genuine mayoral spending preferences}
    \item Municipalities with \textbf{above average fiscal discretion} allocate proportionally \textbf{less} to urbanismo --- when mayors have more fiscal freedom, infrastructure competes less successfully for discretionary funds as other categories absorb the additional room to maneuver
\end{itemize}
\end{frame}


% --- FINDING 2B: FISCAL DISCRETION SLIDE 2 ---%
\begin{frame}{Finding 2: Fiscal Discretion and Spending Composition}
\begin{block}{Result}
The discretionary share of total expenditure ($\tilde{D}_{it}$) significantly shapes spending composition in three categories, independently of electoral incentives
\end{block}
\vspace{1em}
\begin{itemize}
    \item \textbf{Robustness check:} Fiscal discretion is \textbf{not simply a proxy for municipal wealth} --- the correlation between $\tilde{D}_{it}$ and PIB per capita is only 0.12, and with IDHM only 0.21
\end{itemize}
\end{frame}

% --- FINDING 3: INTERACTION ---%
\begin{frame}{Finding 3: The Moderation Effect}
\begin{block}{Result}
The interaction $R_{it} \times \tilde{D}_{it}$ reveals \textbf{heterogeneity in the form of electoral manipulation} depending on the level of fiscal discretion
\end{block}
\vspace{0.5em}
\begin{table}
\centering
\resizebox{\textwidth}{!}{%
\begin{tabular}{lcccc}
\hline\hline
Category & Estimate & Std. Error & T-stat & P-value \\
\hline
Health         & $+$0.079$^{*}$   & 0.047 & 1.68 & 0.094 \\
Social Assist. & $-$0.040$^{***}$ & 0.014 & $-$2.84 & 0.005 \\
Administration & $-$0.149$^{**}$  & 0.057 & $-$2.60 & 0.009 \\
\hline\hline
\multicolumn{5}{l}{\footnotesize $^{***}p<0.01$, $^{**}p<0.05$, $^{*}p<0.10$.}
\end{tabular}%
}
\end{table}
\end{frame}

%--- FINDING 3: INTERPRETATION ---%
\begin{frame}{Finding 3: Conditional Clientelism}
\begin{block}{Unified Interpretation}
The form of electoral manipulation depends on \textbf{how much fiscal room a mayor has}
\end{block}
\vspace{0.5em}
\begin{itemize}
    \item \textbf{Mayors with below average fiscal discretion} rely on \textbf{two electoral instruments simultaneously}:
    \begin{itemize}
        \item \textbf{Patronage hiring} --- expanding the municipal payroll generates electoral loyalty when discretionary funds are scarce, with the effect growing stronger as discretion falls below the mean
        \item \textbf{Social assistance} --- targeting direct transfers to voters becomes the preferred electoral tool when fiscal room is tight
    \end{itemize}
    \vspace{0.5em}
    \item \textbf{Mayors with above average fiscal discretion} show \textbf{attenuated effects} in both categories --- at one standard deviation above the mean, the administration effect ($+1.027$ pp) is nearly fully offset ($-1.12$ pp), suggesting that patronage hiring becomes unnecessary when mayors have ample budget room
\end{itemize}
\end{frame}

\begin{frame}{Finding 3: Conditional Clientelism}
\begin{block}{Unified Interpretation}
The form of electoral manipulation depends on \textbf{how much fiscal room a mayor has}
\end{block}
\vspace{0.5em}
\begin{itemize}
    \item This is consistent with a \textbf{conditional clientelism} story: fiscal constraint shapes the instruments of electoral manipulation, with cash-strapped mayors relying more heavily on patronage and direct transfers
\end{itemize}
\end{frame}

%--- NULL RESULTS ---%
\begin{frame}{The Null Results: Health and Education}
\begin{itemize}
    \item Reelection eligibility has \textbf{no significant effect} on discretionary spending shares for health or education
    \vspace{0.5em}
    \item This is consistent with \textbf{constitutional earmarking} constraining mayoral choices even within the discretionary budget:
    \begin{itemize}
        \item FUNDEB mandates minimum education spending --- mayors face political and legal costs from deviating
        \item SUS transfers create institutional pressures to maintain health spending levels
    \end{itemize}
    \vspace{0.5em}
    \item The null results are a \textbf{substantive finding}, not a failure of the model --- they confirm that institutional constraints on spending are binding even at the level of discretionary allocation
    \vspace{0.5em}
    \item This contributes to the literature on \textbf{fiscal federalism and political accountability}: earmarking effectively insulates key social spending from electoral manipulation, even when mayors have discretionary resources available
\end{itemize}
\end{frame}
\end{document}